% Generated by Listas/make_statements.py from pratica-lecture-01.md, item 23. Do not edit.
\textbf{Source.} Mas-Colell, Whinston \& Green (1995), \emph{Microeconomic Theory}, Ch. 8 Exercises, Exercise 8.B.3\^{}B, p.~262 {[}PDF 142{]}. Reproduced verbatim. Answer in the source: none in the book (see the MWG solutions manual in \texttt{Livros/}).



\begin{quote}

Consider the following auction (known as a \emph{second-price}, or \emph{Vickrey}, auction). An object is auctioned off to \(I\) bidders. Bidder \(i\)'s valuation of the object (in monetary terms) is \(v_i\). The auction rules are that each player submit a bid (a nonnegative number) in a sealed envelope. The envelopes are then opened, and the bidder who has submitted the highest bid gets the object but pays the auctioneer the amount of the \emph{second-highest} bid. If more than one bidder submits the highest bid, each gets the object with equal probability. Show that submitting a bid of \(v_i\) with certainty is a weakly dominant strategy for bidder \(i\). Also argue that this is bidder \(i\)'s unique weakly dominant strategy.

\end{quote}

